\section*{Materials and Methods}

\subsection*{Design, protocol and reporting}
This retrospective evaluation followed a machine-readable protocol registered
on 18 July 2026, before any MIMIC data were accessed, with five dated
amendments before any external result existed. External point estimates were
first inspected on 7 August 2026. Analysis decisions taken after that point,
and every analysis added in a revision audit on 6--7 October 2026, are labelled
post hoc (\ssec{ssec:protocol}). Original estimates are retained beside
corrected ones. Reporting follows TRIPOD+AI where applicable.\cite{collins2024}

This is a secondary analysis of existing de-identified data; no patients were
contacted and no data were collected for it. MIMIC-CXR and MIMIC-CXR-JPG were
accessed through PhysioNet credentialing, which required completion of
human-subjects research training and signature of the PhysioNet Credentialed
Health Data Use Agreement. The Beth Israel Deaconess Medical Center
Institutional Review Board approved the creation of MIMIC-CXR and waived
individual patient consent.\cite{johnson2019cxr} CheXpert Plus was obtained
from the Stanford Center for Artificial Intelligence in Medicine and Imaging
under its research use agreement. No additional institutional ethics review was
sought.

\subsection*{Data, cohorts and context}
The source was MIMIC-CXR-JPG v2.1.0 with free-text reports from
MIMIC-CXR.\cite{johnson2019jpgpaper,johnson2024jpgdata,johnson2019cxr,johnson2024cxrdata,pollard2026}
The external site was CheXpert Plus,\cite{chambon2024} used for evaluation only.
\Cref{tab:cohort} gives the cohort construction. Source studies were split by
patient into model-training, calibration (9,233 studies) and evaluation
(14,766 studies; 5,000 patients) tiers.

Permitted context was the INDICATION and HISTORY report sections at the source
and the published clinical-history fields externally. Extraction was
header-driven: 7.70\% of source reports contain a preliminary interpretation
(WET READ) that a position-based rule would ingest. No extracted context
contained a report-section header. Context named a target finding, usually as a
suspected or known condition, in \ScanTargetSrc{} of source and
\ScanTargetExt{} of external studies. A study was informative (N1) when its
context had at least three alphanumeric tokens after removing
de-identification placeholders. Interventions applied to N1 studies: C0
unchanged; C1 context replaced by a placeholder; C2 context replaced by another
patient's, by deterministic rotation within split and length strata. The audit
found that C2 was applied per image, not per study, so some multi-image
studies received context from several donors (\ssec{ssec:c2}).

\begin{table}[H]
\centering\footnotesize
\caption{Cohort construction (source exclusions are sequential; tiers are
patient-disjoint).}
\label{tab:cohort}
\begin{adjustbox}{max width=\linewidth}
\begin{tabular}{lrr}
\toprule
Step (source site) & Removed & Remaining\\
\midrule
All reports/studies & -- & 227,835\\
No recognised section header & 66 & 227,769\\
Findings and impression merged & 133 & 227,636\\
No frontal (AP/PA) image & 9,688 & 217,948\\
Official split assigned & 0 & 217,948\\
Context not informative ($<3$ effective tokens) & 15,001 & 202,947\\
No impression section & 29,375 & 173,572\\
\midrule
Tier & Patients / studies / frontal images & \\
Model training & 49,959 / 146,253 / 162,823 & \\
Threshold calibration & 3,000 / 9,233 / 10,264 & \\
Prespecified evaluation (primary) & 5,000 / 14,766 / 16,456 & \\
Official validate (not analysed) & 448 / 1,391 / 1,584 & \\
Official test (not analysed) & 282 / 1,929 / 2,155 & \\
\midrule
External site (CheXpert Plus) & & \\
Usable frontal images / studies / patients & 191,070 / 187,673 / 64,709 & \\
Intervention cohort (N1 at $T=3$): images / studies / patients & 135,561 / 132,923 / 53,110 & \\
\bottomrule
\end{tabular}

\end{adjustbox}
\end{table}

\subsection*{Labels}
Targets were cardiomegaly, edema, consolidation, atelectasis and pleural
effusion. Source labels came from a port of CheXbert\cite{smit2020} with a
pinned checkpoint, applied separately to impression (primary) and findings
(registered sensitivity) sections. External labels are the CheXbert labels
distributed with CheXpert Plus; their generation settings are undocumented.
Positive and negative cells were supervised. Uncertain and unmentioned cells
were masked and never mapped to negative. A source study entered a label
cohort only if it had that report section, whereas external studies were
retained whatever sections they had.

\subsection*{Models and frozen policy}
M1 is DenseNet-121 (ImageNet-initialised, $224\times224$); M2 is
BERT-base-uncased (128 tokens). M3 averages M1 and M2 probabilities. M4 fuses
pooled image and text features through a two-layer perceptron, trained end to
end. Models were trained on the model-training tier with masked binary
cross-entropy over labelled cells (AdamW; batch 32; early stopping on
patient-disjoint internal validation). Study probabilities are means over
frontal images. Per-pathology thresholds maximise balanced accuracy on
calibration-tier labelled cells. Study confidence is
$g_i=\min_j|2p_{ij}-1|$, and the abstention cutoff is the 20th percentile of
$g$ on the calibration tier, targeting 80\% coverage. Thresholds and cutoffs
were frozen and applied unchanged at both evaluation sites.

\subsection*{Endpoint and the estimator defect}
For study $i$, let $n_i$ be its number of labelled cells, $w_i$ the number of
wrong decisions among them, and $a_i$ its acceptance. The registered endpoint is
the mean study-level Hamming error among accepted studies. The original
implementation computed $R_{\text{orig}}=\sum_i a_i w_i/\max(n_i,1)/\sum_i a_i$,
which counts accepted studies with $n_i=0$ as correct. The code's own
documentation states the opposite intent. We treat this as a verified defect
and report the corrected estimator, which takes the mean over studies whose
error is defined:
\begin{equation}
R_{\text{eval}}=\frac{\sum_i a_i\,\mathbb{1}[n_i>0]\,w_i/n_i}{\sum_i a_i\,\mathbb{1}[n_i>0]},
\qquad R_{\text{orig}}=R_{\text{eval}}\cdot\Pr(n_i>0\mid a_i=1).
\end{equation}
The original estimator therefore scales error by the evaluable share of
accepted studies, and any cross-site difference in that share appears as a
difference in risk. $R_{\text{eval}}$ is the minimal correction. It applies the
registered definition to studies whose error is defined and changes nothing
else, and the identity above can be checked from published aggregates.
Alternative aggregations are reported as sensitivity analyses rather than
substitutes. $R_{\text{eval}}$ is an error rate over labelled cells, not
a clinical five-finding error rate. Coverage is reported over all studies.

\subsection*{Hypotheses and statistical analysis}
With $R^{s}_{m}(c)$ the selective risk at site $s$ for model $m$ under
condition $c$, the registered hypotheses were as follows.
\begin{itemize}
  \item H1 (confirmatory, M3 and M4): the transfer gap
        $R^{E}_{m}(\mathrm{C0})-R^{S}_{m}(\mathrm{C0})>0$.
  \item H2 (confirmatory): the interaction
        $[R^{E}(\mathrm{C1})-R^{E}(\mathrm{C0})]-[R^{S}(\mathrm{C1})-R^{S}(\mathrm{C0})]>0$.
  \item H3 (confirmatory): H2 for each multimodal model minus H2 for M1.
  \item H4 (secondary): $R_{m}(\mathrm{C2})-R_{m}(\mathrm{C1})\ge0$ within each
        site.
\end{itemize}
Registered sensitivity analyses were coverage targets of 70\% and 90\%,
informativeness thresholds of $T=5$ and $T=10$, and findings-derived labels.
H1 and H2 are confirmed when the 95\% interval lies above zero and the
estimate is at least 0.02. The 0.02 materiality was registered without clinical
justification and is a convention. H4 has no materiality: its pre-results
implementation requires only an interval above zero. M1 ignores text, so the
registered H3 equals H2. The H3 reported earlier, the M3 or M4 transfer gap
minus M1's, was defined after inspection and is reported as post hoc.

Intervals are 95\% percentile intervals from 2,000 patient-clustered bootstrap
replicates (seed 20260718), paired within site and independent between sites.
They condition on the fitted models, the realised calibration sample and the
C2 realisation. Post hoc analyses include:
\begin{itemize}
  \item alternative estimators: cell-pooled, pathology-macro and
        pathology-specific with Holm adjustment;
  \item label handling: uncertain counted as positive or as negative, and
        unmentioned counted as negative;
  \item a controlled label-source decomposition;
  \item per-pathology operating characteristics;
  \item failure-detection AUROC;
  \item comparator selective scores fitted on the source calibration tier only
        (temperature-scaled margin, decision-aware expected loss and
        split-conformal singleton acceptance);
  \item calibration-sample resampling;
  \item extra C2 permutations;
  \item a second training seed for all four models;
  \item view and external demographic strata.
\end{itemize}
Nothing was fitted on external data. Every originally reported estimate was
reproduced exactly from cached predictions, and every frozen threshold from fresh
calibration-tier inference (\ssec{ssec:repro}).
